\documentclass[12 pt]{article}  % sets the font to 12 pt and says this is an 
                                % article (as opposed to book or other documents)
\usepackage{amsfonts, amssymb}  % math fonts and symbols
\usepackage{amsmath}            % math
\usepackage{enumitem}           % change enumeration bullet styles
\usepackage{graphicx}           % graphics
\usepackage{parskip}
\usepackage{xcolor}             % color
\usepackage{hyperref}           % linkable references
\hypersetup{
    colorlinks=true,
    linkcolor=black,
    filecolor=black,      
    urlcolor=black,
    citecolor=black
    }
\usepackage[nottoc,numbib]{tocbibind}
  
%\usepackage{setspace}          % Together with \doublespacing below allows for
                                % doublespacing of the document

\oddsidemargin=-0.5cm
\setlength{\textwidth}{7in}
\addtolength{\voffset}{-20pt}
\addtolength{\headsep}{25pt}
\renewcommand\thefootnote{[\roman{footnote}]} % Bracketed footnotes

                                        % Macro definitions for the header
\newcommand{\myname}{Danny Nguyen}      % Name
\newcommand{\myid}{OU ID: 113687989}    % Student ID
\newcommand{\mysemester}{Summer 2026}   % Semester
\newcommand{\myclass}{DSA 5900}         % Class
\newcommand{\myassignment}{Mid-Semester Progress Report}  % Assignment designator

% Page headers automatically horizontal fill; Name ID Date Class Assignment Page.No.
\pagestyle{myheadings}
\markright{\myname \hfill \myid \hfill \today \hfill \myclass \hfill \myassignment \hfill }

\newcommand{\eqn}[0]{\begin{array}{rcl}}    % begin an aligned equation - allows
                                            % for aligning = or inequalities.
                                            % Always use with $$ $$
\newcommand{\eqnend}[0]{\end{array} }  	    % end the aligned equation

\newcommand{\qed}[0]{$\square$}     % make an unfilled square the
                                    % default for ending a proof

% Common number sets, must be in math mode
\newcommand{\Naturals}[0]{\mathbb{N}}   % Natural numbers
\newcommand{\Integers}[0]{\mathbb{Z}}   % Integers
\newcommand{\Rationals}[0]{\mathbb{Q}}  % Rationals
\newcommand{\Reals}[0]{\mathbb{R}}      % Reals
\newcommand{\Complex}[0]{\mathbb{C}}    % Complex

\newcommand{\heading}[1]{\noindent\large{\textbf{#1}}}  % for sections
\newcommand{\problem}[1]{\noindent\paragraph{#1:}}      % for problem labels

%\doublespacing         % Together with the package setspace above allows for
                        % doublespacing of the document

% Use these commands for self-assessments
\newcommand{\correct}[1]{\colorbox{green}{\checkmark #1}}   % mark correct
\newcommand{\correction}[1]{{\color{red}\textbf{#1}}}       % write a correction

\title{\myassignment \\ \large \myclass--\mysemester \\ \myname}

\begin{document}

\begin{titlepage}
    \begin{center}
        \vfill
        \Huge \myassignment \\
        \vspace*{0.5in}
        \Large \mysemester -- \myclass -- 1 Credit Hour\\
        \vspace*{0.5in}
        \myname
        \vfill
        \begin{tabular}{rl}
            Faculty Supervisor: & Dr. Matthew Beattie \\
            Thesis Advisor:     & Dr. Charles Nicholson
        \end{tabular}
        \vfill
    \end{center}
\end{titlepage}
\pagenumbering{roman}
\tableofcontents
\listoftables
\newpage
\pagenumbering{arabic}

\section{Introduction}
% A brief description of the practicum project context.
% a) Discuss why this project is of interest and of value.
% b) What problems or questions does it help address?
% c) Include a discussion of the research questions/objectives of the sponsor (if applicable)
% and how this project will contribute to those goals.
% d) Use appropriate citations and proper formatting

In the modern media landscape, the velocity and volume at which information 
spreads make manual verification insufficient, leaving professional journalists 
and publishers working under severe time constraints. This problem has 
catalyzed a growing interest in fact-checking systems capable of parsing 
arbitrary text, retrieving evidence, and predicting the veracity of assertions. 
\cite{cheng2025empoweringllmslogicalreasoning}
\cite{guo2022surveyautomatedfactchecking}
However, contemporary semantic pipelines frequently rely on generative large 
language models (LLMs) to output structured configurations like JSON, 
introducing critical flaws: they are highly prone to hallucination, 
fragile under slight parsing variances, and computationally expensive at scale.
\cite{jin2022logicalfallacydetection}
\cite{lalwani2025autoformalizingnaturallanguagefirstorder}
\cite{yang2023couplinglargelanguagemodels}

To overcome these limitations, this practicum aims to define a low-compute, 
symbolic framework to extract the underlying predicate logic from synthetic 
natural language reasoning examples using frozen transformer activations. By 
deploying a compact, open-weights model---Meta's Llama 3.2-3B---this practicum 
extracts high-dimensional activation vectors from the model's intermediate 
residual stream layers. Rather than relying on text generation, these numerical 
parameters will feed into shallow, linear classification probes trained to 
deterministically isolate atomic assertions and identify propositional 
relationships (e.g., implication, negation, and conjunction) directly from the 
embeddings within a model's latent space.
\cite{alain2018understandingintermediatelayersusing} 
\cite{marks2024geometrytruthemergentlinear}

\newpage
\section{Objectives}
% A clear and concise statement of the goals of the practicum project
% a) include the technical project objectives
% b) include individual learning objectives
The overarching goal of this practicum is to develop a robust, parameter-level 
extraction pipeline that bypasses generative LLM bottlenecks to map natural 
language into verifiable logical graphs. The specific project and individual 
objectives are outlined below:

\subsection{Technical Project Objectives}
\begin{itemize}
    \item \textbf{Latent Space Probing Architecture:} Design and train 
    low-compute, linear classification heads (probes) to extract propositional 
    claims and logical relationships directly from the intermediate residual 
    stream of Meta's Llama 3.2-3B model.
    \item \textbf{Symbolic Pipeline Integration:} Construct a deterministic 
    pipeline that translates extracted latent representations into discrete 
    symbolic tokens without relying on stochastic text-generation or JSON 
    parsing.
    \item \textbf{Graph-Based Fact Verification:} Implement a visualization and 
    verification layer that structures the extracted symbols into a formal 
    logical graph, for human interpretation of the underlying logical 
    structures of natural language.
    \item \textbf{Performance Optimization:} Evaluate the classifier framework 
    to achieve high factual and structural precision while reducing 
    inference-time computational overhead compared to baseline autoregressive 
    generative prompting.
\end{itemize}

\subsection{Individual Learning Objectives}
\begin{itemize}
    \item \textbf{Mechanistic Interpretability Mastery:} Gain deep, practical 
    expertise in analyzing the internal geometry of transformers, specifically 
    concerning layer-wise activation tracking, representation dynamics, and 
    internal knowledge localization.
    \item \textbf{Advanced Representation Engineering:} Master state-of-the-art 
    diagnostic probing and activation patching methodologies using specialized 
    interpretability toolkits.
    \item \textbf{Symbolic Architecture Development:} Develop a rigorous 
    engineering workflow for bridging machine representations (vector 
    embeddings) with classical symbolic paradigms (graph structures).
\end{itemize}

\newpage
\section{Data}

\subsection{Ingestion}
% Describe what you did to acquire the data and where you are storing it.  
% What obstacles have you overcome in getting it?
The dataset selected for this project is \textbf{LogicNLI}
\cite{tian-etal-2021-diagnosing}, an open-source diagnostic benchmark hosted by 
the Hugging Face hub. The dataset in JSON was retrieved directly from the 
authors' GitHub repository.
\footnote{https://github.com/omnilabNLP/LogicNLI/tree/main}
There have been no obstacles in acquiring the dataset, as it is publicly 
available and well-documented thanks to the work of the authors. The dataset is 
stored locally in a structured directory within the project repository, 
ensuring easy access for preprocessing and analysis.

\subsection{Exploration}
% Describe the basic characteristics of the data, including whatever plots you 
% have generated so far.
Initial exploratory data analysis revealed a highly uniform, balanced, and 
synthetically controlled dataset structure across its 20,000 entries. The 
target \texttt{labels} field contains a near-perfect balanced split across its 
classification categories (Entailment, Contradiction, Self-Contradiction, and 
Neutral) eliminating the risk of majority-class bias during probe calibration.

A structural analysis of the text fields highlights the specific syntactic attributes of the dataset:

\begin{itemize}
    \item \textbf{Facts Array:} This field contains flat, atomic assertions of 
    absolute truth (e.g., \texttt{["Alice is in the garden.", 
    "Bob is hungry."]}). These strings exhibit zero logical connectives, 
    maintaining a short average string length of 4 to 6 tokens.
    \item \textbf{Rules Array:} This field holds the propositional conditional 
    formulas (e.g., \texttt{["If Alice is in the garden, then Alice is 
    reading."]} or structures introducing logical equivalence and 
    conjunction (\texttt{"and"})). These strings demonstrate higher token 
    complexities and specific syntactic anchors that dictate the relationship 
    edges.
    \item \textbf{Statements and Labels:} The \texttt{statements} field 
    presents the test query, which correlates cleanly to an explicit truth 
    value defined by the formal deduction of the combined rules and facts.
\end{itemize}

Exploratory token length histograms further confirmed tight distribution peaks, 
ensuring that the model's intermediate hidden layers will not suffer from 
sequence padding distortions during batch processing.

\subsection{Preparation}
% Describe any methods you have done to prepare the data for analysis, 
% including missing data imputation, new feature creation, etc.
Because the linear probes require distinct feature-target pairs to map out 
Llama's internal vector space, substantial preprocessing and feature 
engineering were executed to transform the raw dataset into a token-level 
sequence classification task. No missing data or null fields were detected due 
to the synthetic nature of LogicNLI, bypassing the need for standard 
statistical imputation.

The data preparation workflow consists of two critical pipeline stages:
\begin{enumerate}
    \item \textbf{Rule-Based Target Transformation and Fact-Rule Matching:} A 
    custom semantic parsing script was built to map the structural elements of 
    the environment. Within the \texttt{rules} array, regular expressions and 
    token-dependency parsing isolate the relational connectives, explicitly 
    identifying whether a rule represents an implication ($\rightarrow$) or a 
    conjunction ($\land$). Concurrently, a lookup matrix maps these rules 
    directly back to the atomic strings inside the \texttt{facts} array. This 
    enables the framework to identify exactly which baseline facts are being 
    operated on by any given rule, providing clear structural targets for the 
    relational edge probes.
    \item \textbf{Contradiction-Driven Negation Curation:} Because explicit 
    negation ($\neg$) is frequently represented implicitly across text pairs 
    rather than through simple keyword matching, a target-generation strategy 
    was implemented leveraging the dataset's logical symmetries. The pipeline 
    filters entries where the target \texttt{label} is explicitly marked as a 
    \texttt{"contradiction"}. By aligning the queried \texttt{statement} 
    against its corresponding asset in the \texttt{facts} array under this 
    specific label, the pipeline mathematically isolates the exact token span 
    undergoing negation. This contradiction-driven mapping provides a 
    high-precision, curated subset of explicit positive-to-negative vector 
    transformations.
    \item \textbf{Token Span Realignment:} Because character indices shift once 
    strings are converted into transformer sub-tokens, the Hugging Face 
    tokenizer was configured with the \\\texttt{return\_offsets\_mapping=True} 
    parameter. This maps the character spans generated during the rule-parsing 
    and contradiction-alignment steps directly onto the sequence tensor 
    positions of Meta's Llama 3.2-3B model. This preparation ensures that when 
    forward hooks intercept the hidden states of layers 8 through 16, the 
    mean-pooling functions grab the precise vector slices corresponding to the 
    isolated logical components.
\end{enumerate}

\newpage
\section{Methodology}

\subsection{Techniques}
% What is your intended modelling technique?  This should include a description 
% of whether this is a supervised or unsupervised approach, what algorithms you 
% will use to solve your model, your approach to test and training splits, and 
% importantly, how will you validate your findings?
To extract and verify the logical structures embedded within Meta's Llama 
3.2-3B model, this practicum implements a supervised learning paradigm using 
diagnostic probing---specifically, the linear classification approach described 
by Alain, et al.
\cite{alain2018understandingintermediatelayersusing} 
Rather than treating the transformer as an end-to-end text generation system, 
this approach treats the internal intermediate activations as a 
high-dimensional feature space from which structured semantic primitives can be 
decoded. 

\subsubsection{Algorithmic Formulation and Modeling}
The foundational modeling technique utilizes shallow linear classification 
probes optimized via standard Logistic Regression with $L_2$ regularization. 
Mathematically, given a dataset of hidden state activation vectors 
$X \in \mathbb{R}^{d}$ extracted from a specific intermediate transformer layer 
$l$ (where $d = 2048$ for Llama 3.2-3B), the probe seeks to learn a weight 
matrix $W$ and bias vector $b$ to predict the discrete logical class $Y$ (e.g., 
term boundaries or connective relations):

$$\hat{Y} = \sigma(W \cdot X + b)$$

where $\sigma$ represents the softmax function.

\subsubsection{Train-Test Splits and Validation Strategy}
The experimental validation follows a rigorous cross-validation protocol using 
the pre-partitioned LogicNLI splits (\texttt{train}, \texttt{dev}, and 
\texttt{test}):

\begin{itemize}
    \item \textbf{Calibration Split:} The probes are trained using activation 
    vectors extracted strictly from the \texttt{train} split.
    \item \textbf{Model Selection and Layer Auditing (Dev Split):} The 
    validation framework uses the \texttt{dev} split to audit the model 
    layer-by-layer. Probes will be independently trained and evaluated across 
    layers 8 through 16 to map the trajectory of logical synthesis. The layer 
    yielding the highest F1-score on the dev set will be selected as the 
    optimal extraction layer.
    \item \textbf{Generalization Verification (Test Split):} Final model 
    performance is reported strictly using vectors from the unseen 
    \texttt{test} split. Out-of-sample generalization will be quantified using 
    Precision, Recall, and F1-score relative to the ground-truth target 
    structures.
\end{itemize}

\subsection{Process Validation}
% Describe any collaboration you have had with your sponsor or stakeholders to 
% validate your intended approach.  Discuss why you have chosen the Techniques 
% you describe above.
The selection of a parameter-level probing methodology over standard 
autoregressive token generation was heavily shaped by preliminary research 
consultations with my thesis advisor. The overarching academic goal was to 
ensure this project delivered a high-quality, meaningful contribution to 
mechanistic interpretability by focusing on two core criteria: drastically 
increasing the ease of human verifiability in arbitrary text and optimizing 
performance under strict, low-compute constraints.

During the scoping phase, discussions with my advisor focused on the inherent 
limitations of using standard generative workflows for truth and fact-checking 
tasks. Forcing an LLM to output structural text or a natural language 
explanation often introduces stochastic hallucinations and abstract reasoning 
errors that a human reviewer cannot easily audit without manual 
fact-checking~\cite{guo2022surveyautomatedfactchecking}.
To ensure a higher standard of human verifiability, the choice to deploy 
diagnostic linear probes directly on the model's internal residual stream was 
validated. By extracting the raw latent geometries and mapping them directly to 
mathematical nodes and directed edges, we eliminate the black-box nature of 
text generation. This methodology allows an observer to easily verify 
conclusions based purely on the explicit, visualizable logical consistency of 
the statements made in the text.

Furthermore, this technique directly addresses the research requirement for 
compute efficiency. Rather than relying on massive, closed-source models or 
computationally expensive fine-tuning loops, this approach utilizes a compact 
3B-parameter model and shallow classification heads. This design ensures that 
the entire extraction and verification architecture can perform with high 
structural precision while operating entirely within local, consumer-grade 
hardware constraints (such as an RTX 4070 Ti). This validation process 
guarantees that the methodology yields a reproducible, low-overhead framework 
that successfully bridges deep learning vector spaces with verifiable symbolic 
graph networks.

\newpage
\section{Results}
% This section is optional for this progress report.  It corresponds to the 
% Evaluation phase of the data science method.

TBD. I hope I have enough time try single-layer and multi-layer probes, and a 
couple different open models to see how results change by model size.

\newpage
\section{Deliverables}
% This section is optional for this progress report.  It corresponds to the 
% Deployment phase of the data science method.

TBD

\cleardoublepage
\bibliographystyle{plain}
\bibliography{../refs.bib}
\end{document}
